% =============================================================================
% ch02_bio_basis.tex — Chapter 2: Biological Basis | NavigatorCrew
% Compiler: XeLaTeX | Language: Hebrew primary, English technical terms via \en{}
% =============================================================================

\selectlanguage{hebrew}

\section{הבסיס הביולוגי}
\label{sec:bio_basis}

% ── 2.1 Principles of Bat Biosonar ──────────────────────────────────────────
\subsection{עקרונות הסונאר הביולוגי של עטלפים}
\label{subsec:bat_biosonar}

הסונאר הביולוגי של עטלפים, הידוע בשם אקולוקציה, מהווה השראה מרכזית
לפיתוח מערכות ניווט אוטונומיות לרחפנים. עטלפים ממשפחת ה-\en{Vespertilionidae}
פולטים קריאות אולטרה-סוניות בתדרים הנעים בין 20 ל-120 קילו-הרץ ומנתחים
את ההדים החוזרים כדי לנווט בסביבות מורכבות \cite{Simmons1979BatSonar}.
יכולת זו מרשימה במיוחד לאור העובדה שעטלפים פועלים בתנאי תאורה אפסיים,
תוך התמודדות עם רעשי רקע חזקים והפרעות מסביבתם.

שני מנגנוני עיבוד עיקריים מאפיינים את מערכת האקולוקציה העטלפית.
המנגנון הראשון הוא ניתוח השהיית ההד (\en{Time-of-Flight}), המאפשר
לעטלף לאמוד מרחק לעצמים בדיוק של מילימטרים בודדים. המנגנון השני,
מורכב יותר, הוא ניתוח השינויים בתדר עקב אפקט דופלר, המאפשר זיהוי
תנועה יחסית של טרף או מכשולים \cite{GriffinBatEcholocation}.
עטלפים מסוימים אף משתמשים באפנון תדר (\en{FM}) לצד קריאות בתדר קבוע
(\en{CF}), ובכך משיגים איזון בין דיוק בטווח לבין רגישות למהירות.

\begin{figure}[H]
    \centering
    \includegraphics[width=0.92\columnwidth]{figures/fig_bat_echolocation_spectrogram.png}
    \caption{ספקטרוגרם של קריאת אקולוקציה אופיינית מעטלף \textit{Eptesicus fuscus}.
             נראה אפנון התדר היורד (\en{FM sweep}) מ-80 קילו-הרץ ל-25 קילו-הרץ
             לאורך 2.5 אלפיות שנייה.}
    \label{fig:bat_echolocation_spectrogram}
\end{figure}

% ── 2.2 Acoustic Signal Processing ──────────────────────────────────────────
\subsection{עיבוד אותות אקוסטי במערכת העצבים של העטלף}
\label{subsec:bat_neural}

מערכת העצבים השמיעתית של העטלף מבצעת עיבוד מקבילי של מידע אקוסטי
בשני מסלולים נפרדים. המסלול הראשון, המכונה \en{AM}
(\en{Amplitude Modulation}), מתמחה בזיהוי עוצמת ההד ובחילוץ מידע
על גודל ומרקם המטרה. המסלול השני, המכונה \en{FM}
(\en{Frequency Modulation}), מתמחה בזיהוי שינויי תדר עדינים
ובחילוץ מידע על מהירות וכיוון התנועה \cite{Schnitzler1968DSC}.
הפרדה זו בין מסלולי עיבוד מקבילה לעקרונות של עיבוד רב-מודאלי
במערכות רובוטיות מודרניות.

מחקרים אלקטרופיזיולוגיים הראו כי נוירונים בגרעין ה-\en{IC}
(\en{Inferior Colliculus}) של העטלף מגיבים באופן סלקטיבי לשילובים
ספציפיים של תדר והשהייה \cite{Schuller1974DSC}. סלקטיביות זו
מאפשרת יצירת מפה תלת-ממדית של הסביבה בזמן אמת, תוך התעלמות
מהדים חוזרים שאינם רלוונטיים. יכולת הסינון הסלקטיבי הזו רלוונטית
במיוחד למערכות \en{LiDAR} ו-\en{SLAM}, שבהן יש צורך לסנן
רעשי רקע ומדידות חריגות.

\begin{table}[H]
\centering
\begin{tabular}{lcc}
\toprule
מסלול עצבי & פרמטר מחולץ & מקבילה מלאכותית \\
\midrule
\en{AM pathway} & עוצמה, מרקם & עוצמת החזר \en{LiDAR} \\
\en{FM pathway} & מהירות, כיוון & אפקט דופלר ב-\en{FMCW} \\
\en{IC combination} & מיקום תלת-ממדי & \en{Factor Graph SLAM} \\
\bottomrule
\end{tabular}
\caption{השוואה בין מסלולי עיבוד עצביים בעטלף לבין מקבילותיהם המלאכותיות.}
\label{tab:bat_neural_pathways}
\end{table}

% ── 2.3 Evolutionary Adaptation ─────────────────────────────────────────────
\subsection{התאמה אבולוציונית לתנאי סביבה משתנים}
\label{subsec:bat_adaptation}

אחד ההיבטים המרתקים של האקולוקציה העטלפית הוא היכולת להתאים
את פרמטרי הקריאה לתנאי הסביבה המשתנים. עטלפים מגדילים את תדר
החזרה על קריאות (\en{Pulse Repetition Rate}) כאשר הם מתקרבים
למטרה, תופעה המכונה "שלב ההאצה" (\en{Buzz Phase}). במצב זה,
תדר הקריאות עולה מ-10-20 הרץ ל-100-200 הרץ, ומאפשר רזולוציה
מרחבית גבוהה במיוחד \cite{MossEcholocation}.

התאמה דינמית זו של פרמטרי החישה רלוונטית ישירות לתכנון מערכות
ניווט לרחפנים. במקום להשתמש בתדר דגימה קבוע, ניתן לאמץ גישה
אדפטיבית שבה תדר הסריקה של ה-\en{LiDAR} או המצלמה עולה כאשר
הרחפן מתקרב למכשולים או נכנס לסביבה צפופה. גישה זו חוסכת
באנרגיה ומשאבי חישוב בסביבות פתוחות, תוך שמירה על רמת דיוק
גבוהה בסביבות מורכבות \cite{GriffithBatEcholocation}.

נגדיר את תדר החזרה על קריאות כפונקציה של המרחק למכשול הקרוב ביותר
$d_{\min}$:

\begin{equation}
    f_{\text{PRR}}(d_{\min}) = f_{\min} +
    \frac{f_{\max} - f_{\min}}{1 + e^{-\alpha(d_{\min} - d_0)}}
    \label{eq:bat_pulse_rate}
\end{equation}

כאשר $f_{\min} = 10$ הרץ הוא תדר החזרה המינימלי,
$f_{\max} = 200$ הרץ הוא התדר המקסימלי בשלב ההאצה,
$d_0 = 2$ מטר הוא מרחק הסף, ו-$\alpha = 1.5$ מטר$^{-1}$ הוא
מקדם החדות. פונקציה זו מדגימה כיצד תדר הדגימה עולה באופן
הדרגתי ככל שהרחפן מתקרב למכשול.

\begin{equation}
    \Delta t_{\text{min}} = \frac{2 d_{\min}}{c}
    \label{eq:min_time_flight}
\end{equation}

משוואה~\ref{eq:min_time_flight} קובעת את מרווח הזמן המינימלי
הנדרש בין קריאות עוקבות, כדי למנוע חפיפה בין ההד החוזר
לקריאה הבאה. עבור $d_{\min} = 1$ מטר, מתקבל
$\Delta t_{\text{min}} \approx 5.8$ אלפיות שנייה,
המתאים לתדר מקסימלי של כ-170 הרץ.

% ── 2.4 Summary ─────────────────────────────────────────────────────────────
\subsection{סיכום והשלכות לתכנון המערכת}
\label{subsec:bio_summary}

העקרונות הביולוגיים שנסקרו בפרק זה מספקים תשתית רעיונית לתכנון
מערכת הניווט הביומימטית. שלושה עקרונות מרכזיים ייושמו ב-\en{NavigatorCrew}:

\begin{enumerate}
    \item \textbf{עיבוד מקבילי רב-מסלולי:} בדומה למסלולי ה-\en{AM} וה-\en{FM}
    במערכת העצבים של העטלף, המערכת תעבד במקביל מידע מ-\en{LiDAR},
    \en{IMU} ומצלמות.
    \item \textbf{סינון סלקטיבי:} מנגנון הסלקטיביות של ה-\en{IC} יתורגם
    לפונקציות עלות \en{Robust} המפחיתות השפעת מדידות חריגות.
    \item \textbf{התאמה דינמית:} תדר הדגימה ישתנה בהתאם למרחק למכשולים,
    בדומה לשינוי תדר החזרה על קריאות בעטלפים.
\end{enumerate}

פרק זה הניח את היסודות הביולוגיים. בפרק הבא נסקור את מודאליות
החיישנים המלאכותיים המממשות עקרונות אלה.